% ======================================================================
% Ontologie der Kontinua -- Core 1.3
% Cross-K-Modul: crossk_k11_k12.tex
% Cross-Level-Relationen zwischen K11 und minimal spezifiziertem K12
% ======================================================================

\subsubsection{K11–K12-Querverbindung}
\label{sec:crossk-k11-k12}

Der Übergang $K_{11} \rightarrow K_{12}$ ist der terminale Cross-Level-Schritt in
Core~1.3. Während $K_{11}$ ein trans-metatheoretisches Kontinuum mit
universellen Operatoren bietet, die auf ganzen Klassen von Meta-Frameworks wirken,
bindet $K_{12}$ diese oberen Strukturen in die finale wissenschaftliche
Übergreifungs-Hülle ein, in der quersystemische Kompatibilität und globale
Integrierbarkeit bewertet werden.

Die Stufe bleibt abstrakt, ist jedoch kein Platzhalter mehr. Ihre Rolle ist,
die Admissibilitätsbedingungen festzulegen, unter denen die geförderte Hierarchie
als ein kohärenter wissenschaftlicher Stapel behandelt werden kann.

% ----------------------------------------------------------------------
\subsubsection{1. Stufen und ihre Rollen}

\paragraph{K11 (trans-metatheoretisches Kontinuum).}
$K_{11}$ liefert:
\begin{itemize}
  \item Achsen $A^{11}$, die Beziehungen zwischen Klassen von Meta-Frameworks
        beschreiben,
  \item Potentiale $P^{11}$, die trans-meta Kohärenz ausdrücken,
  \item Schwellen $\Theta^{11}$ für universelle Konsistenz und Kompatibilität,
  \item Flüsse $J^{11}$ für Umstrukturierung von Familien von K_{11}-Objekten,
  \item Zyklen $C^{11}$, die die Evolution von Meta-Strukturen regulieren.
\end{itemize}

\paragraph{K12 (formales oberes Kontinuum).}
In Core~1.3 wird $K_{12}$ über:
\begin{itemize}
  \item die Existenz eines Umgebungsraums $M_{12}$ möglicher Achsen,
  \item die Möglichkeit neuer Unterscheidungen, die in $A^{11}$ nicht darstellbar sind,
  \item abstrakte Schwellen $\Theta^{12}$ für strukturelle Konsistenz,
  \item generische Flüsse $J^{12}$, die auf Klassen von $K_{11}$-Objekten wirken,
  \item potenzielle Zyklen $C^{12}$, welche neue hochrangige Strukturen stabilisieren,
\end{itemize}
definiert.

$K_{12}$ ist kein konkretes Domänenmodell. Es ist die globale Integrabilitätslage,
die den Widerspruch-freien Abschluss der geschlossenen Hierarchie begrenzt.

% ----------------------------------------------------------------------
\subsubsection{2. Vererbte / neue Achsen und Schwellen}

\paragraph{Vererbung.}
Alle Achsen, Potentiale und Schwellen von $K_{11}$ sind in $M_{12}$ einbettbar:
\[
A^{11} \subseteq M_{12}, \qquad
\Theta^{11} \subseteq \Theta^{12}.
\]

\paragraph{Neue Unterscheidungen.}
Falls $K_{12}$ existiert, müssen seine Achsen erfüllen:
\[
A^{12} \in M_{12} \setminus A(K_{11}),
\]
konsistent mit Theorem~4 (neue Achsen können nicht selbst von $K_{11}$ erzeugt werden).

\paragraph{Schwellenausweitung.}
$\Theta^{12}$ kann einführen:
\begin{itemize}
  \item \textbf{erweiterte Kohärenzschwellen} über ganze Ketten
        $K_0\dots K_{11}$,
  \item \textbf{Obergrenzbedingungen} (von $M_{12}$ vorgegebene Limits),
  \item \textbf{Kompatibilitätsschwellen} für Strukturen jenseits der Meta-Ebenen.
\end{itemize}

% ----------------------------------------------------------------------
\subsubsection{3. Cross-Level-Flüsse, Zyklen und Spannungen}

\paragraph{Flüsse.}
Die abstrakte Flussform genügt hier:
\[
J^{12} = (J_{\mathrm{lift3}},\ J_{\mathrm{ultra2}},\
          J_{\mathrm{constraint2}},\ J_{\mathrm{embed}}),
\]
bedeutend, dass $K_{12}$ auf $K_{11}$-Strukturen auf dieselbe Weise wie $K_{11}$
auf $K_{10}$ wirkt, jedoch auf höherer Ordnung.

\paragraph{Zyklen.}
Potenzielle $C^{12}$ Zyklen verallgemeinern trans-meta Zyklen:
\[
C_{\mathrm{limit}}:\
\text{K11-Struktur} \rightarrow \text{Lift} \rightarrow
\text{Constraint} \rightarrow \text{Verfeinerung} \rightarrow \text{K11-Struktur}.
\]

\paragraph{Spannung.}
Cross-Level-Spannung folgt dem üblichen Muster:
\[
T_{11,12} =
f(\Theta^{12} - P^{12},\ A^{12} - A^{11},\
  J^{12} - J^{11},\ \Theta^{11} - \Theta^{12}).
\]
Übermäßige Spannung kollabiert $K_{12}$ in $K_{11}$.

% ----------------------------------------------------------------------
\subsubsection{4. Geburts- und Todesbedingungen über die Stufen}

\paragraph{Geburt von $K_{12}$.}
In der geförderten Formalisierung:
\[
T_{11} > \Theta_{11,\mathrm{dim}}
\quad\text{und}\quad
A^{12} \in M_{12},
\]
also müssen neue Unterscheidungen im Raum $M_{12}$ existieren, die die
Repräsentationsfähigkeit von $K_{11}$ überschreiten.

Erweiterung des Zustandsraums:
\[
\Omega(K_{12}) = \Omega(K_{11}) \cup \Delta\Omega_{\mathrm{limit}}.
\]

\paragraph{Todesausbreitung.}
Wenn $\Theta^{11}$ versagt, kann keine $K_{12}$-Struktur gebildet werden.

Wenn $\Theta^{12}$ versagt (erweiterte Inkompatibilität oder Paradox), dann:
\[
\Omega(K_{12}) \rightarrow \varnothing,
\]
während $K_{11}$ intakt bleibt.

% ----------------------------------------------------------------------
\subsubsection{5. Konzeptionelle Beispiele}

\paragraph{Meta-Grenzwert-Erweiterung.}
$K_{12}$ kann Eigenschaften ausdrücken, die sich auf gesamte Ketten von
Abbildungen über Stufen $K_0\dots K_{11}$ beziehen.

\paragraph{Ultrakoherenz.}
Eine universelle Nebenbedingung kann Kohärenz des gesamten Modells unter
Transformationen verlangen, die in $K_{11}$ nicht darstellbar sind.

\paragraph{Kollapsfall.}
Jede Verletzung erweiterter Kohärenz:
\[
T_{11,12} > \Theta^{12}_{\mathrm{coh}},
\]
kollabiert das hypothetische $K_{12}$ zurück in das vollständig spezifizierte
$K_{11}$.

Daher lässt sich der K11→K12-Übergang in Core~1.3 in begrenzter Form als formale
obere Hülle für Quersystem-Kompabilität und globale Integrierbarkeit zusammenfassen,
welche die geförderte Hierarchie wissenschaftlich kohärent hält, während die
Kontinuitäts- und Dimensionsaxiome gewahrt werden.
